\documentclass[12pt]{article}

% ===========================================
%  File: short-note-template.tex
%  Description: Style definitions for Short Note
% ===========================================

\usepackage{ifxetex}
\usepackage{geometry}
\usepackage{xcolor}
\usepackage{titlesec}
\usepackage[fontset=fandol]{ctex}
\usepackage{tikz,eso-pic}
\usetikzlibrary{positioning,calc,shadows,shapes.misc}
\usepackage{fancyhdr}
\usepackage[most]{tcolorbox}
\usepackage{amsmath,amssymb,bm,amsthm}
\usepackage{enumitem}
\usepackage{fontspec}

\geometry{
  paperwidth=7.5in, paperheight=10in,
  left=0.5in, right=0.5in,
  top=0.55in, bottom=0.65in,
  footskip=0.45in
}

\setmainfont{Times New Roman}
\linespread{1.3}
\pagecolor{white}

\definecolor{brand}{HTML}{1A9D8F}
\definecolor{brandD}{HTML}{2A7F6F}
\definecolor{keybg}{HTML}{E8F4F1}
\definecolor{keydark}{HTML}{2F5E58}
\definecolor{secblue}{HTML}{2B44C6}
\definecolor{accentD}{HTML}{D94F4F}
\definecolor{termblue}{HTML}{1A4C9A}
\definecolor{ink}{HTML}{1A1A1A}

\titleformat{\section}[block]{\centering\Large\bfseries\color{keydark}}{}{0em}{}[\vspace{0.2em}\titlerule]
\titlespacing*{\section}{0pt}{1.0em}{0.6em}
\titleformat{\subsection}{\large\bfseries\color{brand}}{}{0em}{$\bullet$\ \ }
\titlespacing*{\subsection}{0pt}{0.8em}{0.3em}

\tcbset{
    elegant_style/.style={
        enhanced, breakable, colback=white, colframe=#1,
        fonttitle=\bfseries, coltitle=white,
        attach boxed title to top left={xshift=12pt, yshift*=-10pt},
        boxed title style={colback=#1, frame hidden, arc=3pt, boxrule=0pt, drop shadow={opacity=0.3, shadow xshift=1pt, shadow yshift=-1pt}},
        boxrule=0.5pt, arc=3pt,
        left=10pt, right=10pt, top=14pt, bottom=10pt,
        drop shadow={opacity=0.06, shadow xshift=2pt, shadow yshift=-2pt},
        before skip=10pt, after skip=10pt
    }
}

\tcbset{
    notion_style/.style={
        enhanced, breakable, frame hidden,
        borderline west={3pt}{0pt}{#1},
        colback=#1!7, coltitle=#1, fonttitle=\bfseries\small,
        attach title to upper={\ \ }, after title={\par\vspace{0.2em}},
        left=8pt, right=8pt, top=6pt, bottom=6pt, arc=2pt,
        before skip=6pt, after skip=6pt
    }
}

\newtcolorbox{KeyBox}{elegant_style=brand, title=Key Idea}
\newtcolorbox{Theorem}[1][]{elegant_style=secblue, title=\if\relax\detokenize{#1}\relax Theorem \else Theorem: #1 \fi}
\newtcolorbox{Definition}[1][]{notion_style=brandD, title=\if\relax\detokenize{#1}\relax Def. \else Def (#1). \fi}
\newtcolorbox{Takeaway}{notion_style=accentD, title=Takeaway}

\newcommand{\padzero}[1]{\ifnum#1<10 0\fi\number#1}
\tcbset{tag_style/.style={on line, boxsep=1pt, left=3pt, right=3pt, top=1pt, bottom=1pt, colback=brand, colframe=brand, fontupper=\bfseries\scriptsize\color{white}, arc=2pt, boxrule=0pt, valign=center, before upper=\strut}}

\newcommand{\MakeShortHeader}[3]{
    \noindent
    \begin{tikzpicture}
        \node[fill=keybg!60, rounded corners=6pt, inner sep=12pt, minimum width=\linewidth, text width=\dimexpr\linewidth-30pt\relax, align=left] (box) {
            {\footnotesize \color{keydark!60} \the\year·\padzero{\the\month}·\padzero{\the\day} \hfill @MIStatlE} \\
            \vspace{0.2em}
            {\fontsize{20pt}{24pt}\selectfont\bfseries\color{keydark} #1} \\[0.2em]
            {\normalsize\color{ink!80} #2}
            \vspace{0.5em}
            \foreach \T in {#3} { \tcbox[tag_style]{\T}\hspace{0.2em} }
        };
        \draw[brandD, line width=3pt] (box.north west) -- (box.south west);
    \end{tikzpicture}
    \vspace{1.0em}
}

\newcommand{\Term}[1]{\textcolor{termblue}{\textbf{#1}}}
\fancyhf{}
\renewcommand{\headrulewidth}{0pt}
\fancyfoot[C]{\footnotesize\color{brand!50} · \thepage \ ·}
\pagestyle{fancy}


\begin{document}

\MakeShortHeader
  {Fisher Information}
  {Regularity, Asymptotic Normality \& Geometry}
  {Statistics, Geometry, MLE}

\begin{KeyBox}
Fisher information is more than a variance lower bound. It acts as the local \Term{geometry} of the parameter space and tells us how sharply the likelihood reacts to perturbations in $\theta$.
\end{KeyBox}

\section{Core Concepts}

\begin{Definition}[Score Function]
The score is the gradient of the log-likelihood:
\[
\mathbf{s}(\theta)=\nabla_{\theta}\log p(x;\theta).
\]
Under standard regularity conditions, $\mathbb{E}[\mathbf{s}(\theta)]=0$.
\end{Definition}

\begin{Theorem}[Cram\'er-Rao Lower Bound]
For any unbiased estimator $\hat{\theta}$,
\[
\mathrm{Var}(\hat{\theta})\succeq \mathcal{I}(\theta)^{-1}.
\]
The inequality says the inverse Fisher information controls the best achievable uncertainty.
\end{Theorem}

\begin{Takeaway}
The template is compact because the important sentence, formal definition, theorem, and takeaway each get a visually distinct block without wasting vertical space.
\end{Takeaway}

\newpage

\MakeShortHeader
  {Expectation as a Linear Operator}
  {A one-page refresher for probability notes}
  {Probability, Cheatsheet, Notes}

\begin{KeyBox}
Expectation is linear even when random variables are dependent:
\[
\mathbb{E}[aX+bY+c]=a\mathbb{E}[X]+b\mathbb{E}[Y]+c.
\]
\end{KeyBox}

\section{Useful Blocks}

\begin{Definition}[Variance]
The variance of $X$ is
\[
\mathrm{Var}(X)=\mathbb{E}\left[(X-\mu)^2\right],\qquad \mu=\mathbb{E}[X].
\]
It measures spread around the mean rather than around zero.
\end{Definition}

\begin{Theorem}[Bias-Variance Identity]
For any estimator $\hat{\theta}$,
\[
\mathbb{E}\left[(\hat{\theta}-\theta)^2\right]=\mathrm{Var}(\hat{\theta})+\mathrm{Bias}(\hat{\theta})^2.
\]
\end{Theorem}

\begin{Takeaway}
This page shows how the same visual system works for a dense formula note. The header, card, theorem block, and narrow margins make it easy to read on mobile.
\end{Takeaway}

\newpage

\MakeShortHeader
  {梯度下降的直觉}
  {Compact bilingual note for screen reading}
  {优化, Bilingual, SGD}

\begin{KeyBox}
梯度下降每一步都沿着局部最陡下降方向移动：
\[
\theta_{t+1}=\theta_t-\eta \nabla f(\theta_t).
\]
当学习率合适时，它会稳定逼近局部最优点。
\end{KeyBox}

\section{为什么这个模板适合短内容}

\begin{Definition}[学习率]
学习率 $\eta$ 控制单步更新的幅度。过大时容易震荡，过小时收敛很慢。这个模板用左侧色条块来呈现定义和提醒，阅读负担很低。
\end{Definition}

\begin{Theorem}[Smooth Objective]
若 $f$ 的梯度是 $L$-Lipschitz，且 $\eta < 1/L$，则梯度下降能够保证目标函数值单调下降。这一类结论很适合放在蓝色 theorem 卡片中。
\end{Theorem}

\begin{Takeaway}
第三页展示了中英混排和公式环境。你可以把它当作知识卡片、讲义单页或自媒体配图底稿。
\end{Takeaway}

\end{document}
